\documentclass[11pt,a4paper]{article}
\usepackage[margin=25mm]{geometry}
\usepackage{amsmath,amssymb,amsthm,booktabs}
\newtheorem{lemma}{Lemma}
\title{Unique overgroups from a normal subgroup of index two}
\author{Human Compression project}\date{}
\begin{document}\maketitle
Write $x_L(p,q)$ for the number of ordered pairs $(r,s)$ of a finite
group $L$ with $|r|=p$, $|s|=q$ and $|rs|=2$. We prove that the
two entries of Connor--Leemans, Table~1, at $(8,10)$ and $(8,20)$
are respectively 4460 and 8880. The inputs are the previously
accepted character evaluations and the complete fixed subgroup
roster of $G=O'N$; the missing step is uniqueness of the maximal
overgroup of each nongenerating pair.

For conjugacy classes of subgroups $K,H\leq G$, suppose every
member of the class of $K$ contains $a$ subgroups in the class of $H$.
If their class lengths are $c_K,c_H$, every $H$ has exactly
\[
                         c_Ka/c_H                           \tag{1}
\]
overgroups in the class of $K$. This follows by counting incidences
in both orders. In particular, an equality $c_Ka=c_H$ proves actual
uniqueness, not merely uniqueness of the overgroup's conjugacy class.

\begin{lemma}\label{square}
Let $M$ be a group, let $E$ be a normal subgroup of index two,
and let $Q\neq E$ be another maximal subgroup. If $Q$ has a unique
subgroup $H$ of index two, then $Q\cap E=H$.
\end{lemma}
\begin{proof}
Since $Q$ is maximal and distinct from $E$, it is not contained
in $E$. Its image in $M/E$ is therefore all of $C_2$, and $Q\cap E$
has index two in $Q$. Uniqueness identifies this intersection with $H$.
\end{proof}

For either of the two requested types, the character calculation
gives $x_M(p,q)=0$ for twelve of the thirteen maximal conjugacy
classes of $G$. The sole remaining class is row 7, of order 161280.
A further useful exclusion is that any subgroup class contained in
one of those twelve zero classes has no such pairs either. This
does not require the pair to generate that subgroup.

Every generated subgroup under consideration contains an element
of order eight and has order divisible by five. The two cyclic
order-eight classes are rows 559 and 560. Following their complete
overgroup relation and retaining orders divisible by five gives
\[
\begin{split}
1,7,8,11,12,15,16,18,19,23,24,30,32,33,34,58,59,60,61,\\
102,103,105,106,133,184,186.
\end{split}                                                   \tag{2}
\]
After the zero-overgroup exclusion and removal of $G$ itself, the
same list survives for both types:
\[
                  7,8,30,33,34,58,59,105,106,133,186.          \tag{3}
\]
The certificate verifies the complete relation and every character
zero used in this step. An order test is used only as a necessary
condition, so the list may include groups without one of the exact
required element orders.

Here are all immediate-overgroup classes for the nonmaximal rows
in (3):
\[
\begin{array}{c|l@{\qquad}c|l}
8&7&30&8\\
33&7&34&7\\
58&8,33&59&8,34\\
105&58&106&59\\
133&30&186&133
\end{array}                                                   \tag{4}
\]
For every edge in this display, formula (1), applied to the literal
cover multiplicity and global class lengths, is exactly one. Thus
every indicated immediate overgroup is unique as an actual subgroup.
The respective multiplicities, in the order
\[
(8,7),(30,8),(33,7),(34,7),(58,8),(58,33),
(59,8),(59,34),(105,58),(106,59),(133,30),(186,133),
\]
are $1,42,56,56,56,1,56,1,1,1,6,1$.

Rows 8, 30, 33, 34, 133 and 186 immediately have a unique maximal
overgroup, in row 7, by following their single-parent chains.
It remains to join the two routes above each of rows 58 and 59.
Fix an actual $H$ in row 58, its unique parent $Q$ in row 33,
and the unique maximal overgroup $M$ of $Q$. The roster gives a
unique subgroup $E$ of index two in $M$, in row 8. It also gives
a unique subgroup of index two in $Q$, namely $H$. Since $Q$ is
maximal in $M$ and $Q\neq E$, Lemma~\ref{square} gives $Q\cap E=H$.
Consequently this $E$ is the unique parent of $H$ in row 8. The
unique maximal overgroups reached through $E$ and $Q$ are the same
actual $M$. The complete parent list (4) now proves that $H$ has
only that maximal overgroup. Replace rows 33 and 58 by 34 and 59
to obtain the other case. Finally, rows 105 and 106 inherit
uniqueness through their sole immediate parents 58 and 59.

We have therefore proved that every proper generated subgroup in
(3), and hence every nongenerating pair of either requested type,
has exactly one maximal overgroup. The maximal subgroup in row 7
is self-normalizing in the simple group $G$, and its class length
is 2857239. Thus the number of generating pairs is exactly
\[
 e_G(p,q)=x_G(p,q)-2857239\,x_{H_7}(p,q).                      \tag{5}
\]
Dividing by the freely acting automorphism group of order
921631011840 gives the following normalized pair counts:
\[\begin{array}{c|rrr}
(p,q)&x_G/|\mathrm{Aut}(G)|&2857239x_{H_7}/|\mathrm{Aut}(G)|&e_G/|\mathrm{Aut}(G)|\\\hline
(8,10)&4469&9&4460\\
(8,20)&8898&18&8880\\
\end{array}\]

There is no additional division for duality. Both results agree
with the fixed Table~1 and had remained intervals under the earlier
maximal-subgroup bounds.

The character and subgroup classifications are borrowed inputs.
The contribution is the actual-overgroup argument, especially
Lemma~\ref{square}'s identification of the two routes, and its
application to the two exact counts. No unproved product of
incidence multiplicities along arbitrary cover paths is used.
There is no new finite group computation, chirality assertion,
principal refutation or whole-source completion claim here.

\begin{thebibliography}{9}
\bibitem{CL} T. Connor and D. Leemans, \emph{Algorithmic enumeration
of regular maps}, Ars Math. Contemp. 10 (2016), 211--222.
\bibitem{roster} T. Connor and D. Leemans, \emph{The subgroup lattice
of $O'N$}, fixed 23-page roster, with exact source joins in the certificate.
\end{thebibliography}
\end{document}
